\documentclass[11pt,twocolumn,a4paper]{article}

\usepackage[utf8]{inputenc}
\usepackage[T1]{fontenc}
\usepackage{amsmath,amssymb,amsthm,mathtools}
\usepackage{bm}
\usepackage{graphicx}
\usepackage[margin=1.1in]{geometry}
\usepackage{enumitem}
\usepackage{booktabs}
\usepackage{float}
\usepackage{natbib}
\usepackage{microtype}
\usepackage{lmodern}
\usepackage{listings}
\usepackage{xcolor}
\usepackage[colorlinks=true,linkcolor=blue,citecolor=blue,urlcolor=blue]{hyperref}
\usepackage{cleveref}

\theoremstyle{plain}
\newtheorem{theorem}{Theorem}[section]
\newtheorem{lemma}[theorem]{Lemma}
\newtheorem{proposition}[theorem]{Proposition}
\newtheorem{corollary}[theorem]{Corollary}
\newtheorem{rulename}[theorem]{Rule}

\theoremstyle{definition}
\newtheorem{definition}[theorem]{Definition}
\newtheorem{principle}[theorem]{Principle}
\newtheorem{example}[theorem]{Example}

\theoremstyle{remark}
\newtheorem{remark}[theorem]{Remark}

\crefname{rulename}{Rule}{Rules}
\Crefname{rulename}{Rule}{Rules}
\crefname{principle}{Principle}{Principles}
\Crefname{principle}{Principle}{Principles}

% --- notation -----------------------------------------------------------
\newcommand{\Rpos}{\mathbb{R}_{>0}}
\newcommand{\N}{\mathbb{N}}
\newcommand{\pow}{\pi}
\newcommand{\Lad}{L}
\newcommand{\Rung}{\gamma}
\newcommand{\Src}{\mathrm{Src}}
\newcommand{\Capset}{\mathrm{Cap}}
\newcommand{\Req}{\mathrm{Req}}
\newcommand{\cost}{\mathrm{cost}}
\newcommand{\vans}{\textsf{answer}}
\newcommand{\vempty}{\textsf{empty}}
\newcommand{\vsurf}{\textsf{unsupported}}
\newcommand{\vtime}{\textsf{exhausted}}
\newcommand{\vref}{\textsf{refused}}
\newcommand{\vstarve}{\textsf{starved}}
\newcommand{\abs}[1]{\left|#1\right|}
\newcommand{\hfq}[1]{\texttt{#1}}

\lstset{
  basicstyle=\ttfamily\small,
  keywordstyle=\bfseries,
  commentstyle=\itshape\color{gray!70!black},
  frame=single,
  framesep=4pt,
  rulecolor=\color{gray!45},
  columns=fullflexible,
  showstringspaces=false,
  breaklines=true,
  xleftmargin=0.5em,
  morekeywords={plan,budget,requests,let,from,ask,with,within,else,fail,
    unresolved,when,starved,emit,partial,map,via,then,expect,union,intersect,
    join,on,filter,where,provenance,divergence,as,ladder,over,power,assert},
}

\title{\textbf{Catalytic Ladders in a Federated Query Language:\\
An Integration Report}}

\author{Kundai Farai Sachikonye\\
\small Technical University of Munich\\
\small AIMe Registry for Artificial Intelligence\\
\small \texttt{kundai.sachikonye@wzw.tum.de}}

\date{\today}



\begin{document}
\maketitle

\begin{abstract}
A ladder composes process steps multiplicatively: each rung carries a power
$\pow\in[0,1]$, the fraction of the remaining gap it closes, and a sequence of
rungs closes $1-\prod_i(1-\pow_i)$. This paper asks a question about that
construct which is not whether it is correct but whether it \emph{fits}: can it
be added to a query language built by other hands, for other purposes, without
breaking what is already there?

We add a \hfq{ladder} step to HFQ, a federated query language with a six-verdict
plan semantics, a capability-containment check that runs before any request is
issued, and a water-filling allocator that distributes a request budget across
steps. The integration is reported as an experiment with a pre-registered
falsification condition: the host ships a suite of sixteen checks, and any
regression in that suite counts against the integration. There is none. The
suite reads $16/16$ before the change and $16/16$ after.

Three findings are worth more than the pass count. First, the ladder needed no
change to the capability checker or the allocator, because both already
dispatch on a guard---is this step a source access?---that a local construct
satisfies by not being one. The host had a zero-cost tier and we did not have to
build it. Second, the ladder's refusal reuses the host's existing blame
discipline: a ladder that cannot reach its declared target reports
\vstarve{} with the predecessor field \emph{empty}, so blame terminates at the
ladder rather than falsely accusing a step that answered correctly. Third, the
host already carries a union bound on translation retention,
$\rho\ge1-\sum_i(1-r_i)$, which is the additive relaxation of the ladder's
multiplicative law. Measuring the two on the same numbers shows the additive
form going \emph{vacuous}---falling below zero, hence holding trivially---on
$51.5\%$ of four-stage chains and $99.6\%$ of eight-stage chains, while the
multiplicative form stays informative throughout. The ladder is therefore not
foreign to the host's theory; it is the tight form of a bound the host already
states loosely, and the looseness is quantified here rather than asserted.

We also report a defect the integration produced and the discipline that caught
it. The first parser read the declared target as an extra rung, because the
token \hfq{power} occurs inside \hfq{expect power}: a plan would then compose a
ladder it did not write. The corpus is real---$400$ enzyme records and $944$
reactions fetched from two public services with different schemas---and then
frozen, because the host forbids network access in adapters by construction and
we did not weaken that rule to run our experiment.
\end{abstract}

\tableofcontents

%==============================================================================
\section{Introduction}
\label{sec:intro}

\subsection{The question}

Federating heterogeneous sources is an old problem with a long literature
\citep{halevy2006pdms}; this paper is not about that problem but about a
smaller one. A construct can be correct and still not fit. A formalism proved sound in
isolation may, on contact with a working system, demand a new cost model, a new
failure mode, a change to an existing check, or an exception to a rule the
system relies on. Each of those is a real cost, and none of them shows up in a
proof.

This paper asks whether the ladder fits. The construct is simple: a
\emph{rung} carries a power $\pow\in[0,1]$, the fraction of a remaining gap it
closes; a \emph{ladder} is a sequence of rungs; and the composite power of a
ladder is
\begin{equation}
\pow(\Lad)\;=\;1-\prod_{i=1}^{n}(1-\pow_i).
\label{eq:compose}
\end{equation}
We add it to HFQ \citep{sachikonye2026hfq}, a federated query language we did
not write, and report what happened. The construct is developed in
\citet{sachikonye2026ladder} and given a typed operational semantics in
\citet{sachikonye2026shkladder}; neither is presupposed here beyond
\cref{eq:compose}.

\subsection{Why this host}

HFQ \citep{sachikonye2026hfq} is a good test because it is opinionated. It has
commitments a careless extension would violate:

\begin{enumerate}[leftmargin=2em,itemsep=0.35em]
\item \textbf{Six verdicts, applied in a fixed order.} A step reports
\vans, \vempty, \vsurf, \vtime, \vref{} or \vstarve, and the order of the rules
that assign them is normative rather than an implementation convenience: each
fires only when every earlier cause has been excluded, which is what makes a
verdict \emph{attributive}.
\item \textbf{One payload rule.} Only \vans{} carries a result. A verdict that
carried a partial payload alongside a failure would reintroduce precisely the
ambiguity the verdict set exists to remove.
\item \textbf{A static capability check.} Whether a step's requirements are
contained in its source's declared capabilities is decided before any request
is issued, and the check is counted, because a bound on its cost is a theorem
of the system.
\item \textbf{No network access, by construction.} Every adapter resolves
against a local fixture. This is stated as a deliberate constraint and not a
temporary limitation: the claims are properties of the compiler, and a live
service can neither confirm nor refute them.
\item \textbf{A blame discipline.} A starved step names the predecessor that
starved it, and a walk up those names terminates. A step that blamed a
predecessor which had in fact answered correctly would corrupt the walk.
\end{enumerate}

Any of these could have been the point at which the ladder did not fit.

\subsection{The falsification condition, fixed in advance}

The host ships a validation suite of sixteen checks. We fixed the following
condition before making any change:

\begin{quote}
\emph{If the suite reports fewer than $16/16$ after the ladder is added, the
integration has failed, and the failure is reported as the result.}
\end{quote}

This is worth stating because the alternative---adding a construct, observing
that one's own new tests pass, and reporting success---would not have tested
integration at all. The host's suite was written without knowledge of the
ladder and has no reason to accommodate it.

The suite reads $16/16$ before the change and $16/16$ after
(\cref{sec:results}). We also record what the change touched, because a
non-regressive change that rewrote half the system would be a different result
from one that did not.

\subsection{What we found}

\paragraph{The zero-cost tier already existed.} The ladder is local: it reaches
no source, issues no request, and demands no capability. We expected to have to
teach the host that such a step exists. We did not. The capability checker
guards on \emph{is this a source access}, not on a list of known local kinds,
so a new local kind is exempted without being mentioned; the allocator does the
same. The host had already been forced to distinguish costed from free steps by
its own set operations, and the ladder inherited that distinction
(\cref{sec:cost}).

\paragraph{The refusal reuses the host's blame discipline.} A ladder whose
declared rungs cannot reach its declared target reports \vstarve{} with the
predecessor field \emph{empty}. That is not an omission: it is the host's
existing idiom for a step whose input arrived intact and whose own declaration
was at fault, and it makes the blame walk terminate at the ladder rather than
accuse a step that answered correctly (\cref{sec:refusal}).

\paragraph{The host already had the additive form of our law.} HFQ bounds
translation retention below by $\rho\ge1-\sum_i(1-r_i)$, the union bound.
\Cref{eq:compose} is its multiplicative counterpart. Measuring both on the same
numbers shows the additive form going vacuous---negative, hence trivially
true---on most chains of four stages or more (\cref{sec:bound}). This is a fact
about the host's own theory that the ladder made visible, and we report it as
such rather than as a criticism.

\subsection{What we do not claim}

We do not claim the ladder makes HFQ better at answering queries; it answers a
different question about a plan than the plan's verdicts do. We do not claim
performance results: the ladder is arithmetic on declared numbers and its cost
is not interesting. We do not claim the integration is complete---
\cref{sec:limitations} lists what a full one would require. And we do not claim
the corpus supports any biological conclusion: it is real data used to give the
plans something real to run on, not evidence about metabolism.


\begin{figure*}[!htbp]
    \centering
    \includegraphics[width=\textwidth]{figures/panel_01.png}
    \caption{\textbf{The host's own suite, and the plans it ships, after the
    extension.}
    \textbf{(A)}~The pre-registered falsification condition. The host's suite
    comprises sixteen checks; sixteen pass and none fails, marked by the cross at
    zero. The suite is invoked as a subprocess and its own summary file is read, so
    the figure plots the host's report rather than our restatement of it. Had this
    bar been non-zero the integration would have failed and the failure would have
    been the result.
    \textbf{(B)}~The negative control, by fixture world. All eleven shipped plans
    that contain no ladder execute, across all three of the host's worlds. The
    world is chosen by the host's own selector; an earlier version of this control
    loaded a single world and reported three failures that were properties of the
    control rather than of the extension.
    \textbf{(C)}~Requests issued by each shipped plan. The plans still reach their
    sources and still spend budget: the control is not passing because nothing
    happens.
    \textbf{(D)}~Verdicts across the shipped plans, by step. The surface is
    concentrated on \textsf{answer} with a ridge at \textsf{starved}, which is the
    distribution the host's suite was written to produce and which the extension
    leaves in place.}
    \label{fig:panel1}
\end{figure*}

%==============================================================================
\section{The Host}
\label{sec:host}

We restate only what the extension touches.

\subsection{Plans and steps}

A plan declares a request budget and a sequence of \hfq{let} bindings
terminated by an \hfq{emit}. A step binds a variable to the result of one
operation: a source access (\hfq{from}), a translation (\hfq{map}), or a set
operation (\hfq{union}, \hfq{intersect}, \hfq{join}, \hfq{filter}). Each step
carries the tuple of variables whose values it consumes.

\subsection{The verdicts}

\begin{definition}[Verdict assignment]
\label{def:verdicts}
A step is assigned a verdict by the following rules, applied in order.
\begin{enumerate}[label=(R\arabic*),leftmargin=3.2em,itemsep=0.2em]
\item If some input has a verdict other than \vans, or has \vans{} with a
payload whose retention fell below the step's declared expectation, then
\vstarve; the diagnosis names the offending predecessor.
\item Else if the step's requirements are not contained in its source's
capabilities, then \vsurf; the diagnosis names the difference.
\item Else if the remaining budget is below the source's cost at the input
cardinality, then \vref; the diagnosis names the shortfall.
\item Else issue the request; if it does not complete within the step budget,
then \vtime.
\item Else if the extracted result set is empty, then \vempty.
\item Else \vans, with the result set as payload.
\end{enumerate}
\end{definition}

\begin{remark}[The order is normative]
\label{rem:order}
Each rule fires only when every earlier cause has been excluded. A step
reporting \vempty{} has therefore established that its inputs were adequate,
its capabilities sufficient, its budget available and its request completed---
and hence that the emptiness is a fact about the data rather than about the
machinery. No rule below (R6) can say that. An extension that inserted a new
rule in the middle of this order would change the meaning of every verdict
below it; \cref{sec:where} explains why the ladder does not.
\end{remark}

\subsection{Capability containment}

\begin{definition}[Containment]
\label{def:containment}
A step against source $\Src$ with requirements $\Req$ is \emph{well-capability}
when $\Req\subseteq\Capset(\Src)$. The check is decided from the declarations
alone, in one membership test per required feature, and issues no request under
any outcome.
\end{definition}

\subsection{Effort and the allocator}

The plan budget is measured in \emph{requests}. Each step is assigned an effort
$e_i$, and the allocator maximises total yield subject to
$\sum_i e_i\le B$; the solution is the standard water-filling characterisation
of a concave programme \citep{boyd2004convex}. Steps whose yield is concave are optimised against a shadow
price; steps that are all-or-nothing are charged first at a fixed cost and
removed from the optimisation, since a step function has no interior optimum.

\begin{remark}[The two-tier regime is prior to this paper]
\label{rem:twotier}
That second clause is the host's own, forced by sources that answer a lookup in
one request and answer nothing more thereafter. It has the side effect that the
host already distinguishes steps that consume requests from steps that do not,
and already charges the latter nothing. The ladder occupies that tier without
enlarging it.
\end{remark}


\begin{figure*}[!htbp]
    \centering
    \includegraphics[width=\textwidth]{figures/panel_02.png}
    \caption{\textbf{The ladder occupies a cost tier the host already had.}
    \textbf{(A)}~Budget spent per step in a plan with two source accesses, one set
    operation and one ladder. The two source steps spend one request each; the
    union and the ladder spend nothing. The ladder's bar is not merely small, it is
    zero.
    \textbf{(B)}~Allocated against spent. Every non-source step is allocated the
    nominal unit the host's water-filling allocator assigns to all-or-nothing steps
    and then spends none of it. That asymmetry is the host's existing conservatism,
    visible here on the union step as much as on the ladder, and the ladder adopts
    it rather than introducing a third regime.
    \textbf{(C)}~Cumulative requests along the plan. The trace rises at each source
    step and is flat across the union and the ladder, so adding ladders to a plan
    cannot make it overrun its budget.
    \textbf{(D)}~Requests spent over the number of source steps and the number of
    ladder steps. The surface is a plane rising in one axis and constant in the
    other: cost depends on source accesses alone.}
    \label{fig:panel2}
\end{figure*}

%==============================================================================
\section{The Ladder Step}
\label{sec:where}

\subsection{Syntax}

One production is added.

\begin{definition}[Added grammar]
\label{def:grammar}
\begin{lstlisting}
let x = ladder over Y
        power P [, power P ...]
        [expect power E]
\end{lstlisting}
The step binds \hfq{x}, consumes the value bound to \hfq{Y}, declares a finite
sequence of rung powers each in $[0,1]$, and optionally declares a target the
composite must attain.
\end{definition}

\begin{example}[A plan with a ladder]
\label{ex:plan}
\begin{lstlisting}
plan ladder_demo {
  budget 50 requests

  let acids = from chebi
      ask descendants_of("CHEBI:1")
      within 10

  let L = ladder over acids
      power 0.45, power 0.30, power 0.55
      expect power 0.70

  emit L
}
\end{lstlisting}
The ladder composes to $1-(0.55)(0.70)(0.45)=0.826\,75$, which clears the
declared target, so the step reports \vans. Raising the target to $0.95$ makes
the same plan report \vstarve{} with the shortfall named.
\end{example}

\subsection{Where it goes, and what it does not touch}

\begin{table}[!htbp]
\centering
\caption{The edit. Two files are modified and two that a reader would expect to
need modification do not.}
\label{tab:edit}
\begin{tabular}{llp{6.4cm}}
\toprule
Component & Change & Why \\
\midrule
Parser & added & One dispatcher branch and one parse function; the step
record gains two fields, both defaulted. \\
Executor & added & One dispatch branch and one handler. \\
Capability checker & \textbf{none} & It guards on \emph{is this a source
access}, so a local step is skipped without being named. \\
Allocator & \textbf{none} & Same guard; a ladder is charged the nominal fixed
cost the host already gives its set operations. \\
Verdict set & \textbf{none} & The refusal reuses \vstarve{} and \vempty. \\
Rule order (R1)--(R6) & \textbf{none} & See \cref{prop:noorder}. \\
\bottomrule
\end{tabular}
\end{table}

\begin{proposition}[The ladder does not disturb the rule order]
\label{prop:noorder}
Adding the ladder step leaves the verdict of every step of every plan that
contains no ladder unchanged, and leaves rules (R1)--(R6) unmodified. A ladder
step is assigned a verdict by (R1), (R5) and (R6) only; rules (R2), (R3) and
(R4) cannot fire on it.
\end{proposition}

\begin{proof}
(R1) is applied by the driver to every step before dispatch, so a ladder whose
input starved is starved by the existing rule and the handler is not reached.
(R2) tests $\Req\subseteq\Capset(\Src)$; a ladder has no source and no
requirements, so the test is vacuous and the checker skips it by its existing
guard. (R3) compares the remaining budget against a source cost; there is no
source, and the step spends nothing. (R4) concerns a request that might not
complete; none is issued. (R5) and (R6) apply as usual to the result set the
step passes through. For a plan containing no ladder, no added code is reached
at all: the parser branch is guarded on a keyword the plan does not contain and
the executor branch on a step kind it does not have.
\end{proof}

\begin{remark}[The guard is what made this cheap]
\label{rem:guard}
Both untouched components test whether a step is a source access, and take the
negative branch for everything else. Had either enumerated the local kinds
instead---set operations and translations, say---a new local kind would have
fallen through to the source path and crashed. The distinction between
``anything that is not a source access'' and ``these particular things that are
not source accesses'' is invisible until someone adds a kind, and it is the
whole reason the integration touched two files rather than four.
\end{remark}

\begin{remark}[Where a fallthrough did bite]
\label{rem:fallthrough}
The executor's dispatch is not written that way. Its final line is an
unguarded fallthrough to the source handler, so an unrecognised step kind is
routed to code that immediately asks a registry for a source named
\hfq{None}. Adding an explicit branch is therefore mandatory rather than
stylistic, and omitting it produces a crash rather than a wrong answer---which
is the better failure, and worth recording as the reason we mention it.
\end{remark}

\subsection{Cost}
\label{sec:cost}

\begin{definition}[The ladder is charged nothing]
\label{def:cost}
A ladder step spends zero budget. It is nonetheless \emph{allocated} the
nominal unit that the allocator assigns to every all-or-nothing step.
\end{definition}

\begin{remark}[A reservation that is never drawn]
\label{rem:reserve}
The asymmetry between an allocation of one and a spend of zero is the host's
existing behaviour for set operations, not something the ladder introduces. It
is conservative: the plan reserves slightly more than it needs, so no plan
overruns its budget because of a ladder. We adopt the existing convention
rather than add a third cost regime, on the principle that an extension should
not multiply the categories its host maintains.
\end{remark}


\begin{figure*}[!htbp]
    \centering
    \includegraphics[width=\textwidth]{figures/panel_03.png}
    \caption{\textbf{One composite, two verdicts, and a blame walk that stops where
    it should.}
    \textbf{(A)}~The same plan under two declared targets. Its rungs compose to
    $0.826\,75$ (horizontal line) in both cases; against a target of $0.70$ the
    step is answered and against $0.95$ it is refused. The verdict is a property of
    the declaration, not of the data, and the input step reports
    \textsf{answer} in both runs.
    \textbf{(B)}~The shortfall the diagnosis names: $0.95-0.826\,75=0.123\,25$,
    drawn as the part of the declared target the ladder does not reach. A refusal
    that named no quantity would be less useful than one that names this one.
    \textbf{(C)}~The blame walk. As implemented it returns the ladder alone; had
    the ladder named its input as the culprit it would have returned two steps and
    accused a step that answered correctly. The right-hand bar is the failure mode
    the empty predecessor field exists to prevent, and it is checked rather than
    assumed.
    \textbf{(D)}~The refusal surface: shortfall over composite and declared target.
    It is zero on and below the diagonal---where the ladder reaches its target---
    and rises above it. The two markers are the two runs of panel~A, one on the
    flat region and one on the slope.}
    \label{fig:panel3}
\end{figure*}

%==============================================================================
\section{The Refusal}
\label{sec:refusal}

\begin{definition}[Refusal to climb]
\label{def:refuse}
A ladder step declaring target $E$ whose rungs compose to $\pow(\Lad)<E$
reports \vstarve. The diagnosis names the declared rungs, the observed
composite, the target and the shortfall $E-\pow(\Lad)$, and sets the named
predecessor to \emph{nothing}.
\end{definition}

\begin{remark}[Why the predecessor field is empty]
\label{rem:blame}
The host's blame walk follows named predecessors from a failed step to the step
that caused it, and terminates when a step names none. A ladder that failed its
target was not starved by its input: the input arrived intact and answered
correctly, and the shortfall is in the declaration. Naming the predecessor
would therefore accuse a step that did nothing wrong and would send the walk
one hop too far. The host already faced this exact situation---a translation
that loses too much of an intact input---and answers it the same way. The
ladder reuses the idiom rather than inventing one.
\end{remark}

\begin{theorem}[The refusal is tight]
\label{thm:tight}
A ladder step is refused if and only if executing its declared rungs would fail
to attain the declared target. No ladder that could have reached its target is
refused, and every refused ladder would have fallen short.
\end{theorem}

\begin{proof}
By \cref{eq:compose} the composite $1-\prod_i(1-\pow_i)$ is exactly the
fraction of the initial gap that applying all $n$ rungs closes. The refusal
tests that quantity against the target. The two verdicts therefore coincide by
construction, and the test is decided from the declared numbers without
executing anything.
\end{proof}

\begin{theorem}[The refusal is non-vacuous]
\label{thm:nonvacuous}
There are plans the language admits on which the refusal fires, and plans on
which it declines to fire.
\end{theorem}

\begin{proof}
The plan of \cref{ex:plan} composes to $0.826\,75$; with target $0.70$ it is
answered and with target $0.95$ it is refused. A ladder declaring no rungs is
rejected at parse time. Each is a plan the grammar admits.
\end{proof}

\begin{remark}[Both directions are required]
\label{rem:bothways}
A predicate that refused every plan would classify as little as one that
refused none. \Cref{thm:nonvacuous} therefore exhibits both, and the
measurements of \cref{sec:results} report the counts of each rather than an
agreement rate alone.
\end{remark}


\begin{figure*}[!htbp]
    \centering
    \includegraphics[width=\textwidth]{figures/panel_04.png}
    \caption{\textbf{The host's union bound is the additive relaxation of the
    ladder's law, and it goes vacuous as chains lengthen.}
    \textbf{(A)}~The two forms on identical retentions drawn from $[0.5,0.99]$,
    $4000$ trials per point. The multiplicative form stays near unity; the additive
    form falls, crosses zero between three and four stages, and enters the shaded
    region. A lower bound on a quantity known to lie in $[0,1]$ that has gone
    negative is true and says nothing.
    \textbf{(B)}~The fraction of trials on which the additive form is vacuous:
    zero at two stages, $0.515$ at four, $0.996$ at eight. Vacuity is the common
    case by four stages.
    \textbf{(C)}~The mean gap between the two forms, which grows monotonically from
    $0.44$ to $2.05$ across the range measured.
    \textbf{(D)}~The gap over chain length and retention. It vanishes at high
    retention and short chains---where the union bound is nearly tight---and grows
    in both directions away from that corner. The comparison is not a criticism:
    the additive form is stated as a bound and carries an injectivity hypothesis,
    and what is measured here is how much it gives up.}
    \label{fig:panel4}
\end{figure*}
%==============================================================================
\section{The Bound the Host Already Had}
\label{sec:bound}

HFQ states a lower bound on the retention of a chain of translations,
\begin{equation}
\rho \;\ge\; 1-\sum_{i=1}^{k}\bigl(1-r_i\bigr),
\label{eq:union}
\end{equation}
where $r_i$ is the retention of stage $i$. This is the union bound. Its
multiplicative counterpart is
\begin{equation}
\rho_{\times} \;=\; 1-\prod_{i=1}^{k}\bigl(1-r_i\bigr),
\label{eq:mult}
\end{equation}
which is \cref{eq:compose} with retentions in place of powers.

\begin{proposition}[The additive form is the relaxation]
\label{prop:relax}
For $r_i\in[0,1]$, $1-\sum_i(1-r_i)\le 1-\prod_i(1-r_i)$, with equality only
when at most one $r_i$ differs from $1$.
\end{proposition}

\begin{proof}
Write $q_i=1-r_i\in[0,1]$. The claim is $\prod_i q_i\le\sum_i q_i$, which holds
because each $q_i\le1$ implies $\prod_i q_i\le q_j$ for any $j$, and
$q_j\le\sum_i q_i$ since the terms are non-negative. Equality forces all but
one $q_i$ to vanish.
\end{proof}

\begin{remark}[Where the relaxation becomes vacuous]
\label{rem:vacuous}
\Cref{eq:union} can be negative, and a negative lower bound on a quantity known
to lie in $[0,1]$ says nothing. \Cref{sec:results} measures how often this
happens as a function of chain length; the answer is that it is the common case
by four stages. We report this because it is a property of a bound the host
states and relies on, and because the ladder is the form that does not have it.
\end{remark}

\begin{remark}[What this is not]
\label{rem:notcriticism}
\Cref{eq:union} carries an injectivity hypothesis and is stated by its authors
as a bound, not as an estimate; a bound that is loose is not a bound that is
wrong. The observation here is that the tight form was available and is the
same object as the ladder's composition law, so the two lines of work meet at a
formula neither was written to share.
\end{remark}


\begin{figure*}[!htbp]
    \centering
    \includegraphics[width=\textwidth]{figures/panel_05.png}
    \caption{\textbf{The frozen corpus: two public services with different schemas.}
    \textbf{(A)}~Sizes and coverage. Four hundred KEGG enzyme records, of which
    $372$ carry at least one reaction; $904$ Reactome reactions, of which $353$
    carry at least one catalyst. The catalyst count was zero on the first
    extraction and the discrepancy was found by querying the live service rather
    than by accepting a plausible zero.
    \textbf{(B)}~Reactions per pathway, ranked. The corpus supplies chains of
    genuinely differing length, the longest containing $58$ reactions, which is what
    makes the chain-length axis of panel~6 a property of data rather than of a
    parameter we chose.
    \textbf{(C)}~Reactions per enzyme record in KEGG, on a logarithmic count axis.
    The distribution is heavy-tailed: most records name few reactions and a few name
    many.
    \textbf{(D)}~Inputs, outputs and catalysts per Reactome reaction. The cloud
    sits mostly in the plane of zero catalysts, with a populated slab above it, and
    the two axes of participants are correlated but not equal---the schema-level
    difference from KEGG, which records substrates and products against an EC
    number instead.}
    \label{fig:panel5}
\end{figure*}

%==============================================================================
\section{The Corpus}
\label{sec:corpus}

The plans need something real to run on. We fetched two public services with
different schemas and different identifier spaces---KEGG
\citep{kanehisa2000kegg} and Reactome \citep{milacic2024reactome}---and then
froze the result.

\begin{table}[!htbp]
\centering
\caption{The frozen corpus. Two services, fetched once, different schemas.}
\label{tab:corpus}
\begin{tabular}{lrl}
\toprule
Source & Records & Schema \\
\midrule
KEGG enzymes    & $400$ & EC number $\to$ reactions, substrates, products \\
Reactome        & $904$ & stable id $\to$ inputs, outputs, catalysts \\
\bottomrule
\end{tabular}
\end{table}

Of the $400$ enzyme records, $372$ carry at least one reaction. Of the $904$
reactions, $353$ carry at least one catalyst and they fall in $19$ pathways,
the longest containing $58$ reactions.

\begin{remark}[Why the corpus is frozen, and why that is not a compromise]
\label{rem:frozen}
The host forbids network access in its adapters by construction, stating that
its claims are properties of the compiler and that a live service can neither
confirm nor refute them. We agree with the reasoning and did not weaken the
rule to run an experiment. The data is real and was fetched once; every
measurement then runs against the frozen snapshot, so the results are
reproducible and no claim depends on a service being reachable.
\end{remark}

\begin{remark}[A field that came back empty, and what we did about it]
\label{rem:emptyfield}
The first extraction returned catalysts for \emph{zero} of the $904$
reactions. Rather than record a zero, we queried the live service directly and
found catalyst records present on reactions our extractor had reported as
having none. The cause was in the extractor: the summary response carries the
catalyst name on the activity entry itself and omits the nested entity our
code required. After the fix the count is $353$. We report this because a
plausible zero is the easiest kind of error to keep---it looks like a finding
about the data---and the only thing that caught it was checking an empty result
against its source instead of accepting it.
\end{remark}


\begin{figure*}[!htbp]
    \centering
    \includegraphics[width=\textwidth]{figures/panel_06.png}
    \caption{\textbf{Composition read against the corpus.}
    \textbf{(A)}~Composite power against the number of rungs at three uniform rung
    powers, with vertical lines at the six longest pathway lengths in the corpus.
    At a rung power of $0.05$ even the longest chain in the corpus falls short of
    unity, which is the saturation result placed against real chain lengths rather
    than against an arbitrary axis.
    \textbf{(B)}~Rungs required to reach three targets, as a function of rung
    power, on a logarithmic axis. The count diverges as the rung power falls, and
    the three curves separate sharply as the target approaches unity.
    \textbf{(C)}~Sensitivity against rung power for a single ladder of corpus
    length. Points lie exactly on $P/(1-\pi_j)$, which increases in $\pi_j$:
    control lies at the strongest rung, not the weakest.
    \textbf{(D)}~The saturation surface: composite power over chain length and
    rung power. It rises to unity in the far corner and stays low along the
    low-power edge however long the chain, which is the statement that a long
    sequence of weak steps does not substitute for a short sequence of strong
    ones.}
    \label{fig:panel6}
\end{figure*}


%==============================================================================
\section{Protocol}
\label{sec:protocol}

Eight experiments run, in \texttt{validation/exp\_integration.py}. Every number
below is read from the JSON they write.

The falsification condition for E1 was fixed before any edit was made to the
host. Two of the eight are negative controls, named with the prefix
\textsc{negative control}, following the host's own convention so they can be
harvested by name. A crash is recorded as a failed experiment rather than
aborting the run.

%==============================================================================
\section{Results}
\label{sec:results}

\subsection{The host suite}

\begin{table}[!htbp]
\centering
\caption{The pre-registered falsification condition.}
\label{tab:hostsuite}
\begin{tabular}{lrr}
\toprule
& Before the ladder & After \\
\midrule
Checks passing      & $16/16$ & $16/16$ \\
Exit code           & $0$     & $0$ \\
Network access      & none    & none \\
\bottomrule
\end{tabular}
\end{table}

The suite is invoked as a subprocess and its own summary file is read, so the
number is the host's report and not our restatement of it. The
\emph{network access} row is included because it is a field the host writes
about itself: adding the ladder did not introduce a network dependency, and the
host records that it did not.

\subsection{The construct through the host's machinery}

The plan of \cref{ex:plan} parses, checks and executes. Its ladder composes to
$0.826\,75$, which is $1-(0.55)(0.70)(0.45)$ to the last digit, computed by the
host's executor from the host's parse of the declared powers. The plan issues
one request, all of it for the source step.

The capability check performs \emph{one} membership test for the whole plan.
That is the source step's single requirement; the ladder contributes none, so
the bound the host proves on the cost of its own static check is untouched.

\subsection{Cost}

\begin{table}[!htbp]
\centering
\caption{Budget spent and effort allocated, per step, in a plan with two source
accesses, one set operation and one ladder.}
\label{tab:cost}
\begin{tabular}{lrr}
\toprule
Step & Spent & Allocated \\
\midrule
\hfq{a}\ \ (source)     & $1.0$ & $1.0$ \\
\hfq{b}\ \ (source)     & $1.0$ & $1.0$ \\
\hfq{u}\ \ (union)      & $0.0$ & $1.0$ \\
\hfq{L}\ \ (ladder)     & $0.0$ & $1.0$ \\
\bottomrule
\end{tabular}
\end{table}

The ladder's row is identical to the union's. This is the substance of the
claim that the ladder occupies a tier the host already had: it is not that the
ladder is cheap, but that its cost profile is indistinguishable from that of a
construct the host shipped. The gap between an allocation of $1.0$ and a spend
of $0.0$ is the host's existing conservatism and not something the extension
introduces.

\subsection{The refusal}

Raising the declared target from $0.70$ to $0.95$ on the same plan changes the
ladder's verdict from \vans{} to \vstarve{}, with the shortfall
$0.95-0.826\,75=0.123\,25$ named in the diagnosis. Three properties are
checked, and all three hold:

\begin{itemize}[leftmargin=1.6em,itemsep=0.25em]
\item The input step reports \vans{} in \emph{both} cases. The refusal is a
property of the declaration, not of the data.
\item The diagnosis names no predecessor.
\item The blame walk from the ladder returns the ladder alone, and does not
reach the input.
\end{itemize}

Had the ladder named its input as the culprit, the walk would have returned two
steps and would have accused a step that answered correctly. This is the
failure mode the empty predecessor field exists to prevent, and it is checked
rather than assumed.

\subsection{The bound the host already had}

\begin{table}[!htbp]
\centering
\caption{The host's additive lower bound against the multiplicative law, on
identical retentions drawn uniformly from $[0.5,0.99]$, $4000$ trials per row.
The additive form is \emph{vacuous} when it falls below zero, since it then
bounds a quantity in $[0,1]$ from below by a negative number.}
\label{tab:bound}
\begin{tabular}{rrrrr}
\toprule
Stages $k$ & Mean additive & Mean multiplicative & Mean gap & Fraction vacuous \\
\midrule
$2$ & $+0.4912$ & $0.9343$ & $0.4431$ & $0.000$ \\
$3$ & $+0.2366$ & $0.9836$ & $0.7470$ & $0.173$ \\
$4$ & $-0.0114$ & $0.9959$ & $1.0073$ & $0.515$ \\
$5$ & $-0.2817$ & $0.9989$ & $1.2806$ & $0.810$ \\
$6$ & $-0.5344$ & $0.9997$ & $1.5342$ & $0.942$ \\
$8$ & $-1.0505$ & $1.0000$ & $2.0505$ & $0.996$ \\
\bottomrule
\end{tabular}
\end{table}

Both the gap and the fraction vacuous increase monotonically with chain length.
By four stages the additive bound is negative on more than half the trials, and
by eight stages on almost all of them. The multiplicative form remains in
$[0,1]$ and informative throughout, by \cref{prop:relax}.

We restate \cref{rem:notcriticism}: the additive form is stated as a bound and
carries an injectivity hypothesis, and a loose bound is not a wrong one. What
the measurement shows is how much is given up, and that the tight form is an
object the host's theory and this construct share.

\subsection{The negative controls}

\paragraph{Plans without a ladder still run.} All eleven shipped plans that
contain no ladder parse and execute, across all three of the host's fixture
worlds. This control failed on first run, and the failure was in the control:
it had loaded a single fixture world and three plans belong to others, so it
reported three unknown-source errors that were properties of our test rather
than of the extension. Using the host's own world selector fixed it. We record
this because a control that fails for its own reasons is worse than no control,
and the only reason we did not report three spurious regressions is that we
read the failures instead of the count.

\paragraph{Malformed ladders are refused.} Four malformed declarations---no
rungs, a power above one, a negative power, and no operand---are all parse
errors. A fifth check verifies that in
\hfq{power 0.5, power 0.6 expect power 0.7} the declared rungs are
$(0.5,0.6)$ and the target is $0.7$.

\begin{remark}[The defect that check exists for]
\label{rem:defect}
That last check is present because the first parser failed it. The pattern
matching a rung looks for the token \hfq{power}, which also occurs inside
\hfq{expect power}, so the declared target was appended as a third rung: the
plan composed a ladder nobody had written, and reported a composite for it. No
verdict would have fired, because nothing failed---the plan simply meant
something other than what it said. This is the class of silent failure the
host's verdict rules exist to prevent, arriving inside an extension added to
the system that prevents it. The fix is to excise the expectation clause before
reading the rungs, and the check that would have caught it is now in the suite.
\end{remark}

%==============================================================================
\section{What Integration Cost}
\label{sec:cost-of-integration}

A pass count says the extension did not break anything. It does not say what it
took. We record the size and shape of the change so that a reader can judge
whether ``it fits'' means anything.

\begin{table}[!htbp]
\centering
\caption{The change, by component.}
\label{tab:diff}
\begin{tabular}{lll}
\toprule
Component & Kind of change & Nature \\
\midrule
Parser            & additive & one dispatcher branch, one function, two
                               defaulted record fields \\
Executor          & additive & one dispatch branch, one handler, one
                               defaulted result field \\
Capability check  & none     & existing guard already exempts local steps \\
Allocator         & none     & existing guard already charges local steps
                               nominally \\
Verdict set       & none     & reuses \vstarve{} and \vempty \\
Rule order        & none     & \cref{prop:noorder} \\
Existing tests    & none     & $16/16$ before and after \\
\bottomrule
\end{tabular}
\end{table}

\begin{remark}[The diff]
\label{rem:diff}
Measured against the host's repository the change is $116$ inserted lines
and one deleted line across two files. The deletion is a comment
enumerating the step kinds, replaced by the same comment with one kind
added, so no existing line of executable code was removed or altered. We
give the number because \emph{additive} is a claim a reader should be able
to check rather than take on trust.
\end{remark}

\begin{remark}[What would have counted as not fitting]
\label{rem:notfitting}
Any of the following would have been a negative result and we state them so the
positive one is falsifiable: a new verdict label, since the host argues its six
are exhaustive; a change to the order of (R1)--(R6), since that order is what
makes verdicts attributive; a third cost tier; a modification to the capability
checker, since a bound on its cost is a theorem; or a network dependency, since
the host forbids one by construction. None was needed. Had one been, the
honest report would have been that the construct requires its host to change,
which is a real cost and not a detail.
\end{remark}

\begin{remark}[The host deserves the credit for two of these]
\label{rem:credit}
The capability checker and the allocator required no change because both guard
on \emph{is this a source access} rather than enumerating the local kinds they
knew about. That is a choice their author made before this construct existed,
and it is why adding a new local kind cost nothing. An enumeration would have
been equally correct at the time and would have made this integration a
four-file change with two opportunities to get it wrong.
\end{remark}

%==============================================================================
\section{Limitations}
\label{sec:limitations}

\begin{enumerate}[leftmargin=2em,itemsep=0.4em]

\item \textbf{The powers are declared, not derived.} In this integration a
plan writes its rung powers as literals. Deriving them from the data a plan
retrieves would be a larger change and would raise the question of what the
derived quantity means for a result set, which we have not answered here.

\item \textbf{The corpus is not evidence about biology.} It is real data, used
so that the plans have something real to run on. No measurement here supports
any claim about metabolism, and the ladder powers in our example plans are
illustrative numbers rather than measured ones.

\item \textbf{The union-bound comparison is over synthetic retentions.}
\Cref{tab:bound} draws retentions uniformly from $[0.5,0.99]$. The
qualitative conclusion---the additive form goes vacuous as chains
lengthen---follows from \cref{prop:relax} and does not depend on the
distribution, but the specific fractions do.

\item \textbf{We did not test the ladder against a live federation.} By
construction: the host forbids it. What this paper establishes is a property
of the compiler, which is the only thing the host's design permits it to
establish, and we inherit that boundary rather than argue with it.

\item \textbf{One construct is not a general claim about extensibility.} The
ladder is local, pure and cheap. A construct that needed to issue requests, or
to fail in a way the six verdicts do not describe, would test the host far
harder, and nothing here predicts how that would go.

\item \textbf{The edits live on the host's main branch.} They are additive and
the host's suite passes, but we have not run them past the host's author.

\end{enumerate}

%==============================================================================
\section{Conclusion}
\label{sec:conclusion}

We set out to find whether a composition law proved in isolation survives
contact with a system built by other hands. The answer, for this construct and
this host, is that it does, and the interesting part is why.

The ladder needed no new cost model because the host had already been forced to
distinguish steps that consume requests from steps that do not. It needed no
new failure mode because the host had already met the situation of a step whose
input was intact and whose own declaration was at fault, and had already
decided that such a step should name no predecessor. It needed no change to the
static checker because that checker asks whether a step reaches a source rather
than which local kinds exist. In each case the accommodation was already there,
made for the host's own reasons, and the extension found it rather than
building it.

Two things went wrong and both are in the record. A regular expression read a
declared target as an extra rung, so a plan would have composed a ladder nobody
wrote---a silent failure inside a system designed to make failures loud. And a
negative control reported three regressions that were artefacts of the control
loading the wrong fixtures; reading the failures rather than the count is what
distinguished them from real ones.

The measurement we did not anticipate is the one about the host's own theory.
HFQ bounds retention below by a union bound, which is the additive relaxation
of the ladder's law. On chains of four stages the relaxation is negative more
than half the time, and by eight stages almost always---a lower bound in
$[0,1]$ that has fallen below zero has stopped saying anything. The tight form
was available all along, and it is the same formula the ladder composes with.
Neither line of work was written with the other in view, and they meet at
$1-\prod_i(1-x_i)$.

\bibliographystyle{plainnat}
\bibliography{references}

\end{document}
